\subsection{Data sources}
\label{sec:data}

Every source is keyed on the CVE identifier, so no vulnerability is duplicated;
each additional database contributes enrichment or PoC rows rather than a second
copy of the record.
All sources are snapshot-pinned by git commit or content hash in a manifest, so
the corpus is reproducible.

\paragraph{Identity and enrichment.}
The \href{https://github.com/CVEProject/cvelistV5}{MITRE CVE List}
(\texttt{cvelistV5}) is the canonical record of identity, CNA provenance, and
description.
\href{https://nvd.nist.gov/vuln/data-feeds}{NVD} supplies normalized CWE, CPE,
and CVSS; \href{https://github.com/cisagov/vulnrichment}{CISA Vulnrichment}
backfills recent CWE/CVSS via its ADP container;
\href{https://github.com/cisagov/kev-data}{CISA KEV} marks known-exploited CVEs;
\href{https://www.first.org/epss/}{FIRST EPSS} supplies exploit-likelihood
scores; and the \href{https://cwe.mitre.org/}{MITRE CWE list} provides the
weakness hierarchy used for classification.
\href{https://github.com/github/advisory-database}{GitHub Security Advisories},
\href{https://access.redhat.com/security/data}{Red~Hat}, and
\href{https://security-tracker.debian.org/tracker/}{Debian} add further CWE,
CVSS, severity, and fix-state signals through their CVE aliases.
Because the same CVE can carry different weakness labels from different sources,
enrichment is kept \emph{source-separated}: each CWE or CVSS row is tagged with
its origin rather than merged, which both preserves provenance and lets us
measure cross-source disagreement.

\paragraph{Proof-of-concept payloads.}
Four PoC databases supply real attack strings:
\href{https://gitlab.com/exploit-database/exploitdb}{Exploit-DB}, the
\href{https://github.com/projectdiscovery/nuclei-templates}{Nuclei detection
templates}, the
\href{https://github.com/rapid7/metasploit-framework}{Metasploit Framework}
modules, and the
\href{https://github.com/nomi-sec/PoC-in-GitHub}{nomi-sec PoC-in-GitHub} index.
The first three are cloned locally so that a payload can be extracted from the
exploit file; the last contributes links only.
PoC references are unified in a single table, with a local path where a payload
is extractable.
Because exploit code is dual-use, the corpus stores only links, identifiers, and
locally-extracted payload fragments; no exploit files are redistributed.

\paragraph{Scale.}
The pinned snapshot comprises \textbf{365{,}310} CVE records
(\textbf{347{,}651} in the \texttt{PUBLISHED} state).
Across all categories the PoC databases link to 25{,}041 CVEs (Exploit-DB),
8{,}975 (PoC-in-GitHub), 4{,}256 (Nuclei), and 3{,}106 (Metasploit); their
coverage of the injection subset is reported with the results in
\cref{sec:results}.
A key methodological caveat is that these sources have different, sometimes
mutable, update cadences (NVD in particular re-scores CVEs continuously), so
the analysis fixes a dated snapshot and treats ``field missing in source'' as
distinct from ``property absent in the vulnerability.''
